\documentclass[11pt]{article}
\usepackage[letterpaper,margin=0.72in]{geometry}
\usepackage[T1]{fontenc}
\usepackage[utf8]{inputenc}
\usepackage{lmodern}
\usepackage{microtype}
\usepackage{xcolor}
\usepackage{graphicx}
\usepackage{booktabs}
\usepackage{longtable}
\usepackage{tabularx}
\usepackage{array}
\usepackage{enumitem}
\usepackage[most]{tcolorbox}
\usepackage{fancyhdr}
\usepackage{titlesec}
\usepackage{hyperref}
\usepackage{xurl}
\usepackage{lastpage}
\usepackage{ragged2e}
\usepackage{multicol}
\usepackage{amssymb}
\definecolor{BAOrange}{HTML}{E87521}
\definecolor{BANavy}{HTML}{17324D}
\definecolor{BASlate}{HTML}{5B6573}
\definecolor{BALight}{HTML}{F4F6F8}
\definecolor{BAGold}{HTML}{D7A62A}
\hypersetup{colorlinks=true,linkcolor=BANavy,urlcolor=BAOrange,citecolor=BANavy,pdftitle={Catalog of 50 Business Analysis Techniques}}
\setlength{\parindent}{0pt}
\setlength{\parskip}{4pt}
\setlist{nosep,leftmargin=*}
\renewcommand{\arraystretch}{1.18}
\newcolumntype{Y}{>{\RaggedRight\arraybackslash}X}
\newcolumntype{L}[1]{>{\RaggedRight\arraybackslash}p{#1}}
\pagestyle{fancy}
\setlength{\headheight}{14pt}
\fancyhf{}
\lhead{\textcolor{BANavy}{\small BABOK v3 Study Series}}
\rhead{\textcolor{BASlate}{\small Catalog of 50 Business Analysis Techniques}}
\cfoot{\small \thepage\ / \pageref{LastPage}}
\renewcommand{\headrulewidth}{0.4pt}
\titleformat{\section}{\Large\bfseries\color{BANavy}}{\thesection}{0.7em}{}
\titleformat{\subsection}{\large\bfseries\color{BAOrange}}{\thesubsection}{0.6em}{}
\titleformat{\subsubsection}{\normalsize\bfseries\color{BANavy}}{\thesubsubsection}{0.5em}{}
\newtcolorbox{keybox}[1][]{enhanced,breakable,colback=BALight,colframe=BANavy,boxrule=0.7pt,arc=1.5mm,left=2mm,right=2mm,top=1.5mm,bottom=1.5mm,#1}
\newtcolorbox{practicebox}[1][]{enhanced,breakable,colback=white,colframe=BAOrange,boxrule=0.8pt,arc=1.5mm,left=2mm,right=2mm,top=1.5mm,bottom=1.5mm,#1}
\newtcolorbox{warningbox}[1][]{enhanced,breakable,colback=white,colframe=BAGold,boxrule=0.8pt,arc=1.5mm,left=2mm,right=2mm,top=1.5mm,bottom=1.5mm,#1}
\begin{document}
\begin{titlepage}
\centering
\vspace*{1.15in}
{\Huge\bfseries\color{BANavy} Catalog of 50 Business Analysis Techniques\par}
\vspace{0.35in}
{\Large\color{BAOrange} Numbered reference, purpose, and practical use\par}
\vspace{0.55in}
\rule{0.78\textwidth}{1.2pt}\par
\vspace{0.35in}
{\large Technical reference and study artifact based on the uploaded \emph{A Guide to the Business Analysis Body of Knowledge (BABOK Guide), Version 3}.\par}
\vfill
\begin{keybox}
\textbf{Source discipline.} BABOK terminology and organization are preserved where practical. Explanations are paraphrased for study and engineering use. Any engineering, security, Agile, UML/BPMN, or enterprise-architecture crosswalks are practitioner mappings, not normative BABOK statements.
\end{keybox}
\vspace{0.35in}
{\small Prepared October 2026\par}
\end{titlepage}
\tableofcontents
\newpage
\section{How to Use the Catalog}
The numbering follows Chapter 10 of BABOK v3. The purpose summaries are paraphrased from the source. The ``strong fit'' column is an editorial navigation aid intended to help practitioners select techniques; it is not an official classification.
\begin{longtable}{L{0.06\textwidth} L{0.24\textwidth} L{0.44\textwidth} L{0.18\textwidth}}
\toprule
\textbf{No.} & \textbf{Technique} & \textbf{Purpose summary} & \textbf{Strong fit}\\
\midrule
\endfirsthead
\toprule
\textbf{No.} & \textbf{Technique} & \textbf{Purpose summary} & \textbf{Strong fit}\\
\midrule
\endhead
10.1 & Acceptance and Evaluation Criteria \textcolor{BASlate}{\scriptsize (p. 217)} & Define the conditions a solution must satisfy and the measures used to compare alternatives. & Acceptance, option selection, testability\\
10.2 & Backlog Management \textcolor{BASlate}{\scriptsize (p. 220)} & Record, order, monitor, and refine remaining work items. & Agile planning, prioritization, delivery sequencing\\
10.3 & Balanced Scorecard \textcolor{BASlate}{\scriptsize (p. 223)} & Assess and manage performance using multiple complementary perspectives rather than a single metric. & Strategy execution, performance management\\
10.4 & Benchmarking and Market Analysis \textcolor{BASlate}{\scriptsize (p. 226)} & Compare performance and market conditions to identify improvement opportunities and stakeholder value. & Competitive analysis, target setting\\
10.5 & Brainstorming \textcolor{BASlate}{\scriptsize (p. 227)} & Generate many ideas rapidly, then group and refine them for further analysis. & Discovery, problem solving, ideation\\
10.6 & Business Capability Analysis \textcolor{BASlate}{\scriptsize (p. 230)} & Use capabilities to clarify outcomes, strategic alignment, scope, and prioritization. & Enterprise architecture, transformation planning\\
10.7 & Business Cases \textcolor{BASlate}{\scriptsize (p. 234)} & Justify a course of action by weighing expected benefits against costs, effort, risk, and other considerations. & Investment decisions, funding, option selection\\
10.8 & Business Model Canvas \textcolor{BASlate}{\scriptsize (p. 236)} & Describe how an enterprise creates, delivers, and captures value. & Business model analysis, innovation\\
10.9 & Business Rules Analysis \textcolor{BASlate}{\scriptsize (p. 240)} & Identify, express, validate, refine, and organize rules that govern business behavior and decisions. & Policy logic, decision requirements, compliance\\
10.10 & Collaborative Games \textcolor{BASlate}{\scriptsize (p. 243)} & Use structured group activities to build shared understanding of a problem, need, or solution. & Workshops, alignment, elicitation\\
10.11 & Concept Modelling \textcolor{BASlate}{\scriptsize (p. 245)} & Organize domain concepts and vocabulary so stakeholders share a consistent semantic model. & Domain analysis, terminology, data/process alignment\\
10.12 & Data Dictionary \textcolor{BASlate}{\scriptsize (p. 247)} & Define data elements consistently so stakeholders interpret data in the same way. & Data governance, interfaces, requirements precision\\
10.13 & Data Flow Diagrams \textcolor{BASlate}{\scriptsize (p. 250)} & Show where data originates, how activities transform it, where it is stored, and where it goes. & Data movement, process analysis, interface analysis\\
10.14 & Data Mining \textcolor{BASlate}{\scriptsize (p. 253)} & Discover useful patterns and insights in data to support decisions. & Analytics, opportunity discovery, performance analysis\\
10.15 & Data Modelling \textcolor{BASlate}{\scriptsize (p. 256)} & Represent data entities or classes, their attributes, and their relationships using common semantics. & Information architecture, persistence, API/data design\\
10.16 & Decision Analysis \textcolor{BASlate}{\scriptsize (p. 261)} & Evaluate decision alternatives and their possible outcomes, especially under uncertainty. & Option analysis, risk-aware choices\\
10.17 & Decision Modelling \textcolor{BASlate}{\scriptsize (p. 265)} & Represent the logic used to make repeatable business decisions. & Rules engines, policy decisions, automation\\
10.18 & Document Analysis \textcolor{BASlate}{\scriptsize (p. 269)} & Extract context, requirements, constraints, and other business analysis information from existing materials. & Legacy analysis, policy review, discovery\\
10.19 & Estimation \textcolor{BASlate}{\scriptsize (p. 271)} & Forecast effort, cost, duration, or other quantities needed to evaluate a course of action. & Planning, budgeting, sizing\\
10.20 & Financial Analysis \textcolor{BASlate}{\scriptsize (p. 274)} & Evaluate the financial characteristics of an investment, solution, or solution approach. & ROI, cost-benefit, investment comparison\\
10.21 & Focus Groups \textcolor{BASlate}{\scriptsize (p. 279)} & Elicit reactions, needs, preferences, and opinions through a moderated group discussion. & Product discovery, user research\\
10.22 & Functional Decomposition \textcolor{BASlate}{\scriptsize (p. 283)} & Break complex processes, systems, functions, or deliverables into smaller analyzable parts. & Scope decomposition, architecture, requirements refinement\\
10.23 & Glossary \textcolor{BASlate}{\scriptsize (p. 286)} & Define key domain terms to create a shared vocabulary. & Terminology control, onboarding, requirements clarity\\
10.24 & Interface Analysis \textcolor{BASlate}{\scriptsize (p. 287)} & Analyze information exchanges across solution components or boundaries, including actors, timing, purpose, and content. & API/integration requirements, system boundaries\\
10.25 & Interviews \textcolor{BASlate}{\scriptsize (p. 290)} & Elicit information through structured or semi-structured conversations with individuals or groups. & Discovery, stakeholder needs, relationship building\\
10.26 & Item Tracking \textcolor{BASlate}{\scriptsize (p. 294)} & Record issues and stakeholder concerns, assign ownership, and monitor them to resolution. & Issue management, decisions, action tracking\\
10.27 & Lessons Learned \textcolor{BASlate}{\scriptsize (p. 296)} & Capture successes, failures, improvement opportunities, and recommendations for future work. & Retrospectives, process improvement\\
10.28 & Metrics and Key Performance Indicators (KPIs) \textcolor{BASlate}{\scriptsize (p. 297)} & Define and use measures that indicate performance of solutions, components, or other stakeholder concerns. & Outcome measurement, operational governance\\
10.29 & Mind Mapping \textcolor{BASlate}{\scriptsize (p. 299)} & Visually organize connected ideas and information around a central subject. & Discovery, decomposition, synthesis\\
10.30 & Non-Functional Requirements Analysis \textcolor{BASlate}{\scriptsize (p. 302)} & Define the quality conditions and operating characteristics a solution must meet in addition to functional behavior. & Performance, security, reliability, usability, operability\\
10.31 & Observation \textcolor{BASlate}{\scriptsize (p. 305)} & Learn about work, context, needs, and performance by watching activities as they occur. & Process discovery, usability, operational analysis\\
10.32 & Organizational Modelling \textcolor{BASlate}{\scriptsize (p. 308)} & Represent roles, responsibilities, relationships, and reporting structures in relation to organizational goals. & Operating model, change impact, accountability\\
10.33 & Prioritization \textcolor{BASlate}{\scriptsize (p. 311)} & Provide a structured basis for deciding the relative importance of business analysis information. & Backlog ordering, release planning, trade-offs\\
10.34 & Process Analysis \textcolor{BASlate}{\scriptsize (p. 314)} & Assess a process for effectiveness, efficiency, and opportunities for improvement. & Process improvement, control analysis\\
10.35 & Process Modelling \textcolor{BASlate}{\scriptsize (p. 318)} & Represent how work is performed so it can be understood, analyzed, communicated, and improved. & BPMN/process maps, operating procedures\\
10.36 & Prototyping \textcolor{BASlate}{\scriptsize (p. 323)} & Create iterative representations of a solution to elicit and validate needs, explore design options, and improve user experience. & UX, requirement validation, design discovery\\
10.37 & Reviews \textcolor{BASlate}{\scriptsize (p. 326)} & Evaluate the content and quality of a work product with appropriate stakeholders. & Quality assurance, approval, defect detection\\
10.38 & Risk Analysis and Management \textcolor{BASlate}{\scriptsize (p. 329)} & Identify uncertainty that may reduce value, assess it, and plan or manage responses. & Risk registers, security risk, change risk\\
10.39 & Roles and Permissions Matrix \textcolor{BASlate}{\scriptsize (p. 333)} & Show responsibilities and access/permission relationships to expose gaps and clarify accountability. & Authorization, RACI-like analysis, segregation of duties\\
10.40 & Root Cause Analysis \textcolor{BASlate}{\scriptsize (p. 335)} & Identify and evaluate the underlying causes of a problem rather than only its symptoms. & Defect prevention, incident/problem analysis\\
10.41 & Scope Modelling \textcolor{BASlate}{\scriptsize (p. 338)} & Define boundaries and show which elements are inside, outside, or interacting across them. & System context, project scope, trust boundaries\\
10.42 & Sequence Diagrams \textcolor{BASlate}{\scriptsize (p. 341)} & Model scenario logic by showing the ordered exchanges among participating objects or components. & UML interaction design, API flows, threat analysis\\
10.43 & Stakeholder List, Map, or Personas \textcolor{BASlate}{\scriptsize (p. 344)} & Identify stakeholders and analyze characteristics such as interests, influence, impact, and needs. & Stakeholder strategy, UX, governance\\
10.44 & State Modelling \textcolor{BASlate}{\scriptsize (p. 348)} & Represent the possible states of an entity, transitions between states, and behavior associated with each state. & Workflow logic, lifecycle design, UML state machines\\
10.45 & Survey or Questionnaire \textcolor{BASlate}{\scriptsize (p. 350)} & Collect structured information from a group efficiently and consistently. & Broad elicitation, user research, benchmarking\\
10.46 & SWOT Analysis \textcolor{BASlate}{\scriptsize (p. 353)} & Assess internal strengths and weaknesses alongside external opportunities and threats. & Strategy, option framing, current-state analysis\\
10.47 & Use Cases and Scenarios \textcolor{BASlate}{\scriptsize (p. 356)} & Describe how an actor or system interacts with a solution to achieve a goal. & Functional requirements, testing, UML use cases\\
10.48 & User Stories \textcolor{BASlate}{\scriptsize (p. 359)} & Capture a small unit of desired functionality or quality from a stakeholder-value perspective. & Agile backlog, incremental delivery\\
10.49 & Vendor Assessment \textcolor{BASlate}{\scriptsize (p. 361)} & Evaluate whether a supplier can meet commitments for delivering and sustaining a product or service. & Third-party risk, sourcing, procurement\\
10.50 & Workshops \textcolor{BASlate}{\scriptsize (p. 363)} & Bring stakeholders together in a facilitated session to achieve a defined collaborative goal. & Elicitation, decisions, conflict resolution\\
\bottomrule
\end{longtable}
\section{Technique Selection by Analytical Need}
The same technique can support multiple tasks. The categories below are intentionally overlapping.
\begin{longtable}{L{0.23\textwidth} L{0.66\textwidth}}
\toprule
\textbf{Need} & \textbf{High-value techniques to consider}\\
\midrule
Discovery / elicitation & Interviews; Workshops; Observation; Focus Groups; Survey or Questionnaire; Document Analysis; Collaborative Games; Brainstorming; Prototyping.\\
Scope / decomposition & Scope Modelling; Functional Decomposition; Business Capability Analysis; Mind Mapping; Stakeholder List, Map, or Personas.\\
Process / behavior & Process Analysis; Process Modelling; Use Cases and Scenarios; Sequence Diagrams; State Modelling; Business Rules Analysis; Decision Modelling.\\
Data / information & Data Modelling; Data Dictionary; Data Flow Diagrams; Data Mining; Concept Modelling; Interface Analysis.\\
Decision / prioritization & Decision Analysis; Prioritization; Business Cases; Financial Analysis; Estimation; Benchmarking and Market Analysis; SWOT Analysis.\\
Quality / acceptance & Acceptance and Evaluation Criteria; Reviews; Non-Functional Requirements Analysis; Metrics and KPIs; Prototyping.\\
Risk / governance & Risk Analysis and Management; Roles and Permissions Matrix; Vendor Assessment; Item Tracking; Lessons Learned.\\
Performance / improvement & Balanced Scorecard; Metrics and KPIs; Process Analysis; Root Cause Analysis; Benchmarking and Market Analysis; Lessons Learned.\\
Agile / incremental delivery & Backlog Management; User Stories; Acceptance and Evaluation Criteria; Prioritization; Prototyping; Workshops; Reviews.\\
\bottomrule
\end{longtable}
\section{Technique Combination Patterns}
\begin{itemize}
\item \textbf{Discovery workshop:} Stakeholder Map $\rightarrow$ Brainstorming $\rightarrow$ Process Modelling $\rightarrow$ Business Rules Analysis $\rightarrow$ Prioritization.
\item \textbf{API/integration analysis:} Scope Modelling $\rightarrow$ Interface Analysis $\rightarrow$ Data Dictionary $\rightarrow$ Sequence Diagrams $\rightarrow$ Non-Functional Requirements Analysis.
\item \textbf{Operational problem:} Metrics/KPIs $\rightarrow$ Process Analysis $\rightarrow$ Root Cause Analysis $\rightarrow$ Risk Analysis $\rightarrow$ Business Case.
\item \textbf{Product increment:} User Stories $\rightarrow$ Acceptance Criteria $\rightarrow$ Prototyping $\rightarrow$ Reviews $\rightarrow$ Backlog Management.
\item \textbf{Enterprise transformation:} Business Capability Analysis $\rightarrow$ SWOT $\rightarrow$ Business Case $\rightarrow$ Functional Decomposition $\rightarrow$ Process Modelling $\rightarrow$ Metrics/KPIs.
\end{itemize}

\section*{Primary Source}
\addcontentsline{toc}{section}{Primary Source}
International Institute of Business Analysis (IIBA), \emph{A Guide to the Business Analysis Body of Knowledge (BABOK Guide)}, Version 3, 2015. This document is a study and practitioner synthesis of the user-supplied copy; it is not an official IIBA publication.
\end{document}
